\scanpage{111}
\mspage{60}
La situation $(X,U,T)$ se reconstitue, à isom. unique près, par la connaissance de

1°) les deux catégories $U,T$

2°) Le bifoncteur
\[
\left\{\begin{aligned} &H:U^{\unsure{\circ|o}}\times T\unsure{\relbar|\longrightarrow}\mathrm{Ens}\\ &H(u,t)=\mathrm{Hom}_X(u,t)\end{aligned}\right.
\]

De plus, ces données 1°) 2°) ne sont soumises à aucune restriction. On a, plus exactement

\uline{Proposition 1} Un isomorphisme entre la catégorie $\unsure{\mathscr T|\mathcal T}$ dont les objets sont les triples $(X,U,T)$, avec $X\in\mathrm{Ob}\,\mathrm{Cat}$, $U\in\mathrm{Ouv}(X)$, $T\in\mathrm{Ferm}(X)$, $U$ et $T$ complémentaires l'un de l'autre, les \deleted{flèches étant les morphismes} \inserted{morphismes de}
$(X,U,T)$ dans $(X',U',T')$ étant les foncteurs $f:X\longrightarrow X'$ qui induisent $U\longrightarrow U'$, $T\longrightarrow T'$ (\unsure{d'-à|c-à-d|d'à d} \unsure{$U=f^{-1}(U')$}, $T=f^{-1}(T')$), avec la catégorie \inserted{$\mathscr T'$} des triples $(U,T,H)$, de deux catégories $U,T$ et d'un bifoncteur
\[
H:U^\circ\times T\longrightarrow\mathrm{Ens}
\]
(ou encore un préfaisceau $U^\circ\longrightarrow T^\vee$ sur $U$, donc un objet de $\underline{\mathrm{Hom}}(U^\circ,T^\vee)$), un morphisme de $(U,T,H)$ dans $(U',T',H')$ étant un triple
\[
(f_U,f_T,u)\quad\left|\begin{aligned}&\left.\begin{aligned}f_U &: U\longrightarrow U'\\ f_T &: T\longrightarrow T'\end{aligned}\right|\ \text{des foncteurs}\\ &u:H\longrightarrow H'\circ(f_U^\circ\times f_T)\quad\text{i.e.}\\ &\quad H(u,t)\longrightarrow H'(f_U(u),f_T(t)\\ &\quad\text{un hom. fonctoriel (flèche}\\ &\quad\text{dans }(U\times T^\circ)^\wedge\simeq\underline{\mathrm{Hom}}(U^\circ,T^\vee)\end{aligned}\right.
\]
composition des flèches évidente.

Si $U\mathrel{\lhook\joinrel\relbar}X\mathrel{\relbar\joinrel\rhook}T$ est un triple, \deleted{\unsure{d'une}} et
\[
f_{\unsure{U'|U}}:U'\longrightarrow U,\ f_{\unsure{T'|T}}:T'\longrightarrow T
\]
des flèches de $\mathrm{Cat}$, il y a une façon privilégiée de ``recoller'' $U'$ et $T'$, de façon à ce que le l'objet $(X',U',T')$ s'envoie dans $(X,U,T)$ \deleted{à \unsure{l'aide des}} en induisant
% archive sheet number: 111
